% BHU PYQ -- Philosophy: Indian Methods of Argumentation (2025-26 NEP)
\documentclass[11pt,a4paper]{article}
\usepackage[a4paper,top=2.5cm,bottom=2.5cm,left=2.8cm,right=2.8cm,headheight=14pt]{geometry}
\usepackage{amsmath,amssymb}
\usepackage{fontspec}
\setmainfont{Kohinoor Devanagari}
\usepackage{microtype,fancyhdr,enumitem,hyperref,lastpage,parskip}

\hypersetup{colorlinks=true,urlcolor=black,linkcolor=black,pdftitle={PHISE31: Indian Methods of Argumentation -- BHU NEP 2025-26}}

\pagestyle{fancy}\fancyhf{}
\renewcommand{\headrulewidth}{0.4pt}\renewcommand{\footrulewidth}{0.4pt}
\fancyhead[L]{\small \textit{Banaras Hindu University}}
\fancyhead[R]{\small \textit{PHISE31: Indian Methods of Argumentation (2025-26)}}
\fancyfoot[C]{\small Page \thepage\ of \pageref{LastPage}}

\newlist{parts}{enumerate}{2}
\setlist[parts,1]{label=(\alph*),leftmargin=*,itemsep=4pt,topsep=3pt}
\setlist[parts,2]{label=(\roman*),leftmargin=*,itemsep=2pt,topsep=2pt}

\newcommand{\pts}[1]{\hfill \textit{[#1]}}

\begin{document}

\begin{center}
  {\large\bfseries BANARAS HINDU UNIVERSITY}\\[2pt]
  {\normalsize B. A. (Hons.) Semester III Examination, 2025-26 (NEP)}\\[6pt]
  \rule{0.8\linewidth}{0.5pt}\\[6pt]
  {\large\bfseries PHILOSOPHY}\\[4pt]
  {\large\bfseries Paper Code: PHISE31 : Indian Methods of Argumentation}\\[4pt]
  {\normalsize Time: 2:30 Hours \hfill Full Marks: 70}
\end{center}
\rule{\linewidth}{0.4pt}
\smallskip

\noindent \textit{Instructions: Attempt all three sections A, B and C according to the instructions given. Write your Roll No.\ at the top immediately on receipt of this question paper.}

\smallskip
\noindent \textit{दिये गये निर्देशों के अनुसार खण्ड अ, ब तथा स के उत्तर दीजिये।}

\bigskip

\section*{SECTION -- A / खण्ड -- अ}
\noindent \textit{Note: Attempt any \textbf{three} questions, each in 300 words. Each question carries 10 marks.}\pts{3 $\times$ 10 = 30}

\noindent \textit{किन्हीं तीन प्रश्नों के उत्तर दीजिये, प्रत्येक उत्तर 300 शब्दों में हो। प्रत्येक प्रश्न 10 अंकों का है।}

\begin{enumerate}[label=\textbf{\arabic*.},leftmargin=*,itemsep=12pt]

  \item Explain the Upanisadic method of dialogue with suitable examples.\pts{10}
  
  उपनिषदीय सम्वाद विधि की व्याख्या उपयुक्त उदाहरणों के साथ कीजिये।

  \item How many major methods of argumentation have you learned? Explain the differences among them.\pts{10}
  
  आपने तर्क/वाद-विवाद की कितनी प्रमुख विधियाँ सीखी हैं? उनके मध्य अन्तर की व्याख्या कीजिये।

  \item How were the structures of debate formulated in classical Indian texts? Substantiate your answer with suitable references.\pts{10}
  
  शास्त्रीय भारतीय ग्रन्थों में वाद-विवाद की संरचनाएँ किस प्रकार निर्मित की गईं? उपयुक्त सन्दर्भों के साथ अपने उत्तर की पुष्टि कीजिये।

  \item What is meant by Naya in Jain epistemology? How does it help in determining partial viewpoints of truth?\pts{10}
  
  जैन प्रमाणमीमांसा में नय से क्या अभिप्राय है? यह सत्य के आंशिक दृष्टिकोणों को निर्धारित करने में किस प्रकार सहायता करता है?

  \item Write a short note on the Nyaya method of argumentation.\pts{10}
  
  न्याय दर्शन की तर्क-पद्धति पर एक संक्षिप्त टिप्पणी लिखिये।

\end{enumerate}

\bigskip

\section*{SECTION -- B / खण्ड -- ब}
\noindent \textit{Note: Attempt any \textbf{four} questions, each in 150 words. Each question carries 5 marks.}\pts{4 $\times$ 5 = 20}

\noindent \textit{किन्हीं चार प्रश्नों के उत्तर दीजिये, प्रत्येक उत्तर 150 शब्दों में हो। प्रत्येक प्रश्न 5 अंकों का है।}

\begin{enumerate}[label=\textbf{\arabic*.},leftmargin=*,start=6,itemsep=10pt]

  \item Explain the Dialectical Method in Indian Philosophy.\pts{5}
  
  भारतीय दर्शन में द्वन्द्वात्मक पद्धति की व्याख्या कीजिये।

  \item What is Nigrahasthana and in what situations does it mark the end of an argument?\pts{5}
  
  निग्रहस्थान क्या है और किन परिस्थितियों में यह वाद-विवाद के अन्त को चिह्नित करता है?

  \item What is Syadvada in Jain philosophy? Explain its significance in understanding reality.\pts{5}
  
  जैन दर्शन में स्याद्वाद क्या है? यथार्थ को समझने में इसके महत्त्व की व्याख्या कीजिये।

  \item What is the Samvada Vidhi (method of dialogue) in Buddhism?\pts{5}
  
  बौद्ध दर्शन में सम्वाद विधि क्या है?

  \item Prepare a debate structure on Artificial Intelligence (AI) based on your favorite system of argumentation.\pts{5}
  
  तर्क की अपनी पसंदीदा पद्धति पर आधारित कृत्रिम बुद्धिमत्ता (AI) पर एक वाद-विवाद संरचना तैयार कीजिये।

  \item How is syllogism used as a method of argumentation in Buddhism?\pts{5}
  
  बौद्ध दर्शन में तर्क की एक विधि के रूप में न्याय-वाक्य (अनुमान) का उपयोग किस प्रकार किया जाता है?

\end{enumerate}

\bigskip

\section*{SECTION -- C / खण्ड -- स}
\noindent \textit{Note: Attempt all \textbf{ten} questions of this section in a maximum of two sentences each. Each question carries 2 marks.}\pts{10 $\times$ 2 = 20}

\noindent \textit{इस खण्ड के सभी दस प्रश्नों के उत्तर अधिकतम दो वाक्यों में दीजिये। प्रत्येक प्रश्न 2 अंकों का है।}

\begin{enumerate}[label=\textbf{\arabic*.},leftmargin=*,start=12,itemsep=8pt]
  \item 
  \begin{parts}
    \item What are the three forms of debate mentioned in classical Indian texts? Write their names.\pts{2}
    
    शास्त्रीय भारतीय ग्रन्थों में उल्लिखित कथा (वाद-विवाद) के तीन प्रकार कौन-से हैं? उनके नाम लिखिये।

    \item What is the purpose of Vada according to the Nyaya school?\pts{2}
    
    न्याय सम्प्रदाय के अनुसार वाद का उद्देश्य क्या है?

    \item Which Jain concept explains the multi-sidedness of reality in argumentation?\pts{2}
    
    तर्क में यथार्थ की बहुआयामी प्रकृति को कौन-सी जैन अवधारणा स्पष्ट करती है?

    \item Which method of reasoning is used in the Upanisads for philosophical dialogue?\pts{2}
    
    दार्शनिक सम्वाद के लिए उपनिषदों में तर्क की किस पद्धति का उपयोग किया जाता है?

    \item How many types of Naya are accepted in Jain epistemology? Write their names.\pts{2}
    
    जैन प्रमाणमीमांसा में कितने प्रकार के नय स्वीकार किये गये हैं? उनके नाम लिखिये।

    \item Give one example of Chhala (deceptive argument).\pts{2}
    
    छल (छलपूर्ण तर्क) का एक उदाहरण दीजिये।

    \item Who is a disputant?\pts{2}
    
    वादी/विवादी कौन होता है?

    \item Write the names of the five steps of syllogism in Nyaya (Panca-Avayava).\pts{2}
    
    न्याय के पञ्च-अवयव अनुमान के पाँच चरणों के नाम लिखिये।

    \item Which two steps are accepted during argumentation (Vada) in Buddhism?\pts{2}
    
    बौद्ध दर्शन में वाद-विवाद के दौरान कौन-से दो चरण स्वीकार किये जाते हैं?

    \item Write the name of an Upanisadic dialogue (Samvada).\pts{2}
    
    किसी एक उपनिषदीय सम्वाद का नाम लिखिये।
  \end{parts}
\end{enumerate}

\end{document}
